\DocumentMetadata{
  pdfversion=2.0,
  pdfstandard=ua-2,
  lang=en-US,
  tagging=on,
  tagging-setup={math/setup=mathml-SE,tabsorder=structure}
}
\documentclass[11pt,letterpaper]{article}
\usepackage[margin=0.68in,includefoot]{geometry}
\usepackage{newcomputermodern}
\usepackage{amsmath,amscd,mathtools}
\usepackage{tikz,adjustbox}
\providecommand{\square}{\mdlgwhtsquare}
\usepackage[hidelinks]{hyperref}
\hypersetup{
  pdftitle={Probability PhD Exam, May 2005},
  pdfauthor={University of Florida Department of Mathematics},
  pdfsubject={Graduate mathematics examination},
  pdfdisplaydoctitle=true,
  bookmarksopen=true
}
\setcounter{secnumdepth}{0}
\setlength{\parindent}{0pt}
\setlength{\parskip}{0.28em}
\setlength{\emergencystretch}{3em}
\setlength{\leftmargini}{2.15em}
\setlength{\leftmarginii}{2.65em}
\setlength{\leftmarginiii}{3em}
\pagestyle{empty}
\raggedbottom
\allowdisplaybreaks
\NewTaggingSocketPlug{block/list/label}{examproblem}{%
  \tagstructbegin{tag=Lbl}\tagstructend%
  \tagstructbegin{tag=\LBody}%
  \tagstructbegin{tag=\ProblemHeadingTag,title-o={\ProblemHeadingText}}%
  \tagmcbegin{tag=\ProblemHeadingTag,actualtext={\ProblemHeadingText}}%
  #2%
  \tagmcend%
  \tagstructend%
}
\newenvironment{examproblems}[2]{%
  \def\ProblemHeadingTag{#2}%
  \AssignTaggingSocketPlug{block/list/label}{examproblem}%
  \begin{enumerate}%
  \setcounter{enumi}{#1}%
  \renewcommand{\labelenumi}{\textbf{\theenumi.}}%
  \setlength{\topsep}{0.35em}%
  \setlength{\partopsep}{0pt}%
  \setlength{\itemsep}{0.48em}%
  \setlength{\parsep}{0.18em}%
}{\end{enumerate}}
\newenvironment{examsubparts}{%
  \AssignTaggingSocketPlug{block/list/label}{default}%
  \begin{enumerate}%
  \setlength{\topsep}{0.18em}%
  \setlength{\partopsep}{0pt}%
  \setlength{\itemsep}{0.16em}%
  \setlength{\parsep}{0.1em}%
}{\end{enumerate}}
\newenvironment{examinstructions}{%
  \par\begingroup\itshape
}{%
  \par\endgroup\vspace{0.18em}
}
\makeatletter
\renewcommand\subsection{\@startsection{subsection}{2}{0pt}%
  {-1.25ex plus -.25ex minus -.1ex}%
  {0.55ex plus .1ex}%
  {\normalfont\large\bfseries}}
\renewcommand{\@maketitle}{%
  \newpage\null\vskip -1em%
  {\footnotesize\bfseries
    \href{https://www.ufl.edu/}{University of Florida}, \href{https://math.ufl.edu/}{Department of Mathematics}\par}%
  \vskip -0.5em%
  \tagstructbegin{tag=H1,title={Probability PhD Exam, May 2005}}%
  \tagpdfparaOff%
  \tagmcbegin{tag=H1}%
    {\LARGE\bfseries\href{https://gma.math.ufl.edu/exams/probability-phd/}{Probability PhD Exam}, May 2005\par}%
  \tagmcend%
  \tagpdfparaOn%
  \tagstructend%
  \vskip 0.25em%
  \tagmcbegin{artifact}%
    \hrule height 1pt%
  \tagmcend%
  \vskip 1em%
}
\makeatother
\title{Probability PhD Exam, May 2005}
\author{}
\date{}
\begin{document}
\maketitle
\thispagestyle{empty}
\begin{examproblems}{0}{H2}
\def\ProblemHeadingText{Problem 1.}
\item
Let \(X_1, X_2\) be two independent random variables with the same uniform distribution on \((\theta-1/2,\theta+1/2)\), and let \(Y_1=\min(X_1,X_2)\), \(Y_2=\max(X_1,X_2)\),
\begin{examsubparts}
\item[\textnormal{(a).}]
Find \(\mathbb{P}(Y_1\leq\theta\leq Y_2)\),
\item[\textnormal{(b).}]
Find \(\mathbb{P}(Y_1\leq\theta\leq Y_2\mid Y_2-Y_1\geq 1/2)\).
\end{examsubparts}
\def\ProblemHeadingText{Problem 2.}
\item
Let \(\phi(t)\) be a characteristic function, prove that
\begin{examsubparts}
\item[\textnormal{(a).}]
\(1-\operatorname{Re}(\phi(2t))\leq 4(1-\operatorname{Re}(\phi(t)))\),
\item[\textnormal{(b).}]
\(1-|\phi(2t)|^2\leq 4(1-|\phi(t)|^2)\).
\end{examsubparts}
\def\ProblemHeadingText{Problem 3.}
\item
Let \(\{X_n,n\geq 1\}\) be a sequence of i.i.d. nonnegative random variables. let \(S_0=0\), and \(S_n=X_1+\cdots+X_n\). For \(t>0\), we define
\[
\{\omega\mid N_t(\omega)=n\}=\{\omega\mid S_n(\omega)\leq t<S_{n+1}(\omega)\}
\]
 show that
\[
\mathbb{E}(N_t(\omega))=\sum_{n=1}^{\infty}\mathbb{P}(S_n(\omega)\leq t).
\]
\def\ProblemHeadingText{Problem 4.}
\item
Let \(X_1,X_2,\ldots\) be a sequence of strictly positive random variables such that
\[
\mathbb{E}(X_{n+1}\mid\mathcal{F}_n)=f_n(X_n).
\]
 For \(n\geq 2\), let
\[
M_n=\frac{X_1X_2\cdots X_n}{f_1(X_1)f_2(X_2)\cdots f_{n-1}(X_{n-1})}.
\]
\begin{examsubparts}
\item[\textnormal{(a).}]
Show that for \(n\geq 2\), \(M_n\) is a \(\mathcal{F}_n\)-martingale.
\item[\textnormal{(b).}]
Does this martingale converges almost surely and in \(L^1\)? Explain it.
\end{examsubparts}
\def\ProblemHeadingText{Problem 5.}
\item
Let \(\{M_n,n\geq 0\}\) be a sequence of integrable random variables adapted to a filtration \(\mathcal{F}_n\). Assume that for each bounded stopping time~\(T\), \(\mathbb{E}(M_T)=\mathbb{E}(M_0)\), show that \(\{M_n,n\geq 0\}\) is a martingale.
\def\ProblemHeadingText{Problem 6.}
\item
Let~\(X\) and~\(Y\) be random variables such that \(\mathbb{E}(X^2)<\infty\) and \(\mathbb{E}(Y^2)<\infty\), \(\mathbb{E}(X\mid Y)=Y\) and \(\mathbb{E}(Y\mid X)=X\). Show that \(X=Y\) a.s.
\def\ProblemHeadingText{Problem 7.}
\item
Let \(X,Y\) be two independent random variables with \(\mathbb{E}[Y]=0\). Show that for \(p\geq 1\)
\[
\mathbb{E}[|X|^p]\leq\mathbb{E}[|X+Y|^p].
\]
\def\ProblemHeadingText{Problem 8.}
\item
Let \((X,Y)\) be a random point on a unit circle with uniform distribution, that is
\[
\mathbb{P}((X,Y)\in A)=\frac{\operatorname{length}(A)}{2\pi}
\]
 for any Borel subset~\(A\) of \(C_2=\{(x,y)\mid x^2+y^2=1\}\). Find the marginal distribution of~\(X\).
\def\ProblemHeadingText{Problem 9.}
\item
Let \(\mathcal{F}_n\) be a filtration, \(|X_n|\leq Y\), \(Y\) integrable. Suppose that \(X_n\longrightarrow X\) a.e. Using the martingale convergence theorem to prove that
\[
\mathbb{E}[X_n\mid\mathcal{F}_n]\longrightarrow\mathbb{E}[X\mid\mathcal{F}_\infty]\qquad\text{a.e.}
\]
\def\ProblemHeadingText{Problem 10.}
\item
Let \(Y\in L^p\), \(|X_n|\leq Y\) and \(X_n\longrightarrow X\) in distribution. Show that \(X_n\) converges to~\(X\) in \(L^p\).
\end{examproblems}
\end{document}
